\documentclass[11pt]{article}
\usepackage[T1]{fontenc}
\usepackage[utf8]{inputenc}
\usepackage{lmodern}
\usepackage[ngerman]{babel}
\usepackage{amsmath,amssymb,mathtools,array,booktabs,pdflscape,adjustbox,diagbox}
\usepackage{geometry}
\usepackage{microtype}
\geometry{a4paper,margin=1.45cm}
\setlength{\parindent}{1.25em}
\setlength{\parskip}{0pt}
\makeatletter
\renewcommand\footnotesize{\@setfontsize\footnotesize{7.5}{8.5}\selectfont}
\makeatother
\providecommand{\widebar}[1]{\overline{#1}}
\providecommand{\A}{\Delta}
\providecommand{\D}{\Delta}
\allowdisplaybreaks
\setlength{\emergencystretch}{10em}
\sloppy
\hbadness=10000
\vbadness=10000
\hfuzz=2pt
\vfuzz=10pt
\raggedbottom

% current section packet macros
\providecommand{\Alg}{\mathcal A}
\providecommand{\Abar}{\overline{A}}
\providecommand{\Bbar}{\overline{B}}
\providecommand{\Kbar}{\overline{K}}
\providecommand{\Zbar}{\overline{Z}}
\providecommand{\Gcal}{\mathfrak G}
\providecommand{\Hcal}{\mathfrak H}
\providecommand{\Scal}{\mathfrak S}
\providecommand{\Tcal}{\mathfrak T}
\providecommand{\Zcal}{\mathfrak Z}
\providecommand{\Kcal}{\mathfrak K}
\providecommand{\Pcal}{\mathfrak P}
\providecommand{\rad}{\operatorname{rad}}
\providecommand{\Sp}{\operatorname{Sp}}
\providecommand{\End}{\operatorname{End}}
\providecommand{\Mat}{\operatorname{Mat}}
\providecommand{\Inv}{\operatorname{Inv}}

\begin{document}

\clearpage

\subsection*{§6. Galoissche Theorie einfacher Systeme}

Der verschärfte Vertauschungssatz aus §5, 4. drückt die galoissche Theorie einfacher Systeme in bezug auf das Zentrum als Grundbereich aus. Wir betrachten erst Körper, dann allgemein einfache Systeme.

\paragraph{1. Galoissche Theorie der Körper von endlichem Rang nach dem Zentrum.}
Die galoissche Gruppe \(\Gcal\) von \(A\) ist definiert als Gruppe aller Automorphismen, die Fortsetzung des identischen des Zentrums \(P\) sind, also nach §5, 2. als Gruppe aller inneren Automorphismen. Es wird somit
\[
        \Gcal \simeq A^*/P^*,
\]
wo \(A^*\) beziehungsweise \(P^*\) die multiplikative Gruppe der von Null verschiedenen Elemente aus \(A\) beziehungsweise \(P\) bedeuten. Den Untergruppen \(\mathfrak H\) von \(\Gcal\) entsprechen also eineindeutig die \(P^*\) umfassenden Untergruppen \(H^*\) von \(A^*\), vermöge
\[
        \mathfrak H \sim H^*/P^* .
\]
Eine Untergruppe \(\mathfrak H\) von \(\Gcal\) heißt abgeschlossen, wenn \(H^*\) nach Zufügung der Null in einen Körper \(H\) übergeht. Die Tatsache, daß \(C\) die Gruppe \(\mathfrak H\) gestattet, ist gleichbedeutend damit, daß \(C\) mit \(H\) elementweise vertauschbar ist. Damit geht der Vertauschungssatz §5, 4. über in den

\paragraph{Hauptsatz der galoisschen Theorie nichtkommutativer Körper.}
Die Körper \(C\) zwischen \(P\) und \(A\) und die abgeschlossenen Untergruppen \(\mathfrak H\) von \(\Gcal\) lassen sich eineindeutig einander so zuordnen, daß \(C\) voller Invariantenkörper von \(\mathfrak H\) und \(\mathfrak H\) volle Invariantengruppe von \(C\) wird.

\paragraph{2. Galoissche Theorie der einfachen Systeme.}
Ist \(A_f\) einfaches System mit Zentrum \(P\), so wird die galoissche Gruppe \(\Gcal\) wieder die Gruppe der inneren Automorphismen, also isomorph zu
\[
        A_f^*/P^*,
\]
wo jetzt \(A_f^*\) die multiplikative Gruppe der regulären Elemente aus \(A_f\) bedeutet. Eine Untergruppe \(\mathfrak H\sim H^*/P^*\) von \(\Gcal\) heißt einfach-abgeschlossen, wenn der aus \(H^*\) erzeugte Unterring \(H\) einfaches System ist und wenn \(H^*\) aus der Gesamtheit der regulären Elemente von \(H\) besteht. Der Übergang vom Vertauschungssatz zur galoisschen Theorie ist hier gegeben durch den

\paragraph{Hilfssatz.}
Jedes \(P\) umfassende einfache System aus \(A_f\) wird von der Gruppe \(H^*\) seiner regulären Elemente erzeugt. Das ist klar, wenn \(P\) unendlich viele Elemente besitzt; denn das mit Unbestimmten gebildete „allgemeine Element“ ist regulär, und durch geeignete Spezialisierung lassen sich reguläre Basiselemente nach \(P\) bilden. Der Hilfssatz gilt aber nach Shoda allgemein.\footnote{Shoda, Anm. 3) zitierte Arbeit. Dort wird die Einführung der Unbestimmten vermieden.}

Vermöge des Hilfssatzes ist die Vertauschbarkeit von \(S\) mit einem einfachen System \(\mathfrak S\) wieder gleichwertig damit, daß \(S\) die durch die regulären Elemente von \(\mathfrak S\) induzierten Automorphismen gestattet. Somit geht auch hier der Vertauschungssatz §5, 4. über in den

\paragraph{Hauptsatz der Galoisschen Theorie einfacher Systeme.}
Die einfachen, \(P\) umfassenden Systeme \(S\) aus \(A_f\) und die einfach-abgeschlossenen Untergruppen \(\mathfrak H\) von \(\Gcal\) lassen sich eineindeutig einander so zuordnen, daß \(S\) voller Invariantenbereich von \(\mathfrak H\) und \(\mathfrak H\) volle Invariantengruppe von \(S\) wird.

\paragraph{3. Beweis vermöge des Fortsetzungsprinzips.}
Ich gebe für den Fall eines Körpers noch einen zweiten Beweis des Hauptsatzes, der auf dem Prinzip der Fortsetzung von Isomorphismen, also Darstellungen ersten Grades, beruht und, da er ohne Rangabzählung auskommt, weniger Endlichkeitsvoraussetzungen macht. Insbesondere zeigt sich so, in welcher Richtung eine Übertragung auf Körper \(S\) von unendlichem Rang möglich sein wird. Der Beweis benutzt auch nicht die Tatsache, daß es nur eine irreduzible reziproke Darstellungsklasse gibt, gilt daher genau so im Kommutativen; vgl. §8, 2. Ich schicke einige Hilfssätze voraus.

\paragraph{Voraussetzungen.}
Das einfache System \(S\) über \(P\) lasse eine reziproke Darstellung ersten Grades in \(A\) zu, sei also Körper; dabei kann \(A\) von endlichem oder unendlichem Rang über seinem Zentrum \(P\) sein. Eine Zerlegung von \(S_A\) in \(n\) einfache, operatorisomorphe Rechtsideale nach §4, 3. sei gegeben durch
\[
        S_A=r_1+\cdots+r_n
            =e_1S_A+\cdots+e_nS_A
            =e_1A+\cdots+e_nA ;
\]
die durch \(e_i\) erzeugte reziproke Darstellung ist also definiert durch
\[
        e_i\sigma=e_i\sigma_i .
\]

\paragraph{Hilfssatz 1.}
Die \(n\) durch \(e_1,\ldots,e_n\) erzeugten reziproken Darstellungen sind verschieden, also verschiedene Darstellungen derselben Darstellungsklasse. \(S\) besitzt also mindestens so viele verschiedene Darstellungen, wie sein Rang beträgt.

Zum Beweis erweitern wir die Darstellungen aus \(S\) nach §2 Schluß zu Zuordnungen von \(S_A\), die nach \(A\) operatorhomomorph werden. Diese Zuordnungen sind hier definiert durch
\[
        e_i w=e_i\omega
\]
mit \(\omega\in A\) für jedes \(w\in S_A\). Würde nun durch \(e_i\) und \(e_j\), \(i\ne j\), dieselbe Darstellung erzeugt, so käme also \(e_iw=e_jw\) für jedes \(w\in S_A\). Das ist unmöglich wegen \(e_ie_i=e_i\) und \(e_je_i=0\).

\paragraph{Hilfssatz 2.}
Ist \(T\) Körper zwischen \(P\) und \(S\), ist \(s\) der Rang von \(S\) nach \(T\), so läßt jeder reziproke Isomorphismus von \(T\) in \(A\) mindestens \(s\) verschiedene Fortsetzungen zu.

Der reziproke Isomorphismus, also die reziproke Darstellung von \(T\) in \(A\), sei vermittelt durch eine Zerlegung
\[
        T_A=E_1T_A+\cdots+E_hT_A
            =E_1A+\cdots+E_hA,
        \qquad E_it=E_i\tau ,
\]
wobei die \(E_i\) Idempotente sind. Das ist keine Beschränkung der Allgemeinheit, da eine durch ein Basiselement \(E_i\alpha\) erzeugte Darstellung auch durch \(\alpha^{-1}E_i\alpha\), also ein Idempotent, erzeugt wird. Rechtsmultiplikation mit \(S\) ergibt
\[
        S_A=E_1S_A+\cdots+E_hS_A .
\]
Dabei sind die \(E_iS_A\) alle vom Rang \(s\) nach \(A\). Denn der Rang von \(E_iS_A\) ist höchstens \(s\), wie die Einsetzung einer \(T\)-Basis in \(S\) vermöge der Vertauschbarkeit von \(S\) und \(A\) zeigt:
\[
        S=Tw_1+\cdots+Tw_s,
\]
\[
        E_iS_A=E_iT_Aw_1+\cdots+E_iT_Aw_s
              =E_iw_1A+\cdots+E_iw_sA .
\]
Da die Summe \(S_A\) vom Rang \(n=hs\) ist, hat jedes \(E_iS_A\) genau den Rang \(s\). Somit gilt
\[
        E_iS_A=r_{i1}+\cdots+r_{is}
              =e_{i1}A+\cdots+e_{is}A,
\]
wo die \(s\) durch \(e_{i1},\ldots,e_{is}\) erzeugten Darstellungen nach Hilfssatz 1 alle verschieden werden. Sie sind sämtlich Fortsetzungen der durch \(E_i\) erzeugten Darstellung von \(T\); denn aus \(E_it=E_i\tau\) folgt
\[
        (e_{i1}+\cdots+e_{is})t=(e_{i1}+\cdots+e_{is})\tau
\]
und damit \(e_{ij}t=e_{ij}\tau\) wegen der direkten Summenzerlegung.

\paragraph{Definition.}
Die durch \(T\) induzierte Klasseneinteilung \(\mathfrak I_T\) der reziproken Isomorphismen \(\mathfrak I\) von \(S\) auf \(A\) ist so definiert: es werden alle und nur die Isomorphismen aus \(\mathfrak I\) als äquivalent betrachtet, die auf \(T\) denselben Isomorphismus induzieren.

\paragraph{Hauptsatz.}
Ist \(\mathfrak I_T\) die durch \(T\) induzierte Klasseneinteilung, so ist \(T\) maximaler Unterkörper gegenüber dieser Klasseneinteilung.

Denn ist \(Z\) Zwischenkörper von \(T\) und \(S\), vom Grad \(l\) nach \(T\), so spaltet sich nach Hilfssatz 2 jede Klasse aus \(\mathfrak I_T\) in mindestens \(l\) Klassen aus \(\mathfrak I_Z\) auf. Der Übergang von dieser Fassung des Hauptsatzes zu der üblichen entsteht dadurch, daß die Isomorphismenmenge \(\mathfrak I\) durch Multiplikation ihrer Elemente mit dem Inversen irgendeines festen in die Automorphismengruppe übergeht, wobei eine Klasse aus \(\mathfrak I_T\) in die Invariantengruppe von \(T\) übergeht. Die Aussage, daß die abgeschlossenen Untergruppen volle Invariantengruppen sind, ist darin mitenthalten; man hat nur für \(\mathfrak H\sim H^*/P^*\) den Körper \(H\) als Zwischenkörper \(T\) einzusetzen.

\subsection*{§7. Die Klasse ähnlicher Algebren. Zerfällungskörper}

Von jetzt an wird es sich nur um Körper \(A,\overline A\) und Matrizenringe \(A_f\) handeln, die von endlichem Rang nach ihrem Zentrum \(P\) sind; also um einfache, normale Algebren über \(P\), die kurz als Algebren über \(P\) oder auch nur als Algebren bezeichnet werden sollen. \(A,\overline A\) sind dann Divisionsalgebren. In eine Klasse ähnlicher Algebren \(A,A_f\) werden alle \(A_f\) mit demselben zugeordneten \(A\) zusammengefaßt. Isomorphe \(A\) werden dabei identifiziert.

\paragraph{1. Die Gruppe der Algebrenklassen.}
Sind \(A,B\) Algebren über \(P\), so gilt das gleiche für ihr Produkt, das direkte Produkt nach \(P\); denn nach §3 Schluß kommt
\[
        A_f\times_P B_g=C_h,
\]
wo \(C\) eine bis auf Isomorphie eindeutig bestimmte Divisionsalgebra mit dem Zentrum \(P\) ist. Die Multiplikation von \(A_f,B_g\) mit Matrizenringen über \(P\) zeigt, daß diese Multiplikation eindeutig für die Klassen bestimmt ist, die somit ein multiplikativ abgeschlossenes, kommutatives System bilden. Darüber hinaus gilt der Satz von R. Brauer: die Klassen ähnlicher Algebren bilden gegenüber der direkten Produktbildung eine abelsche Gruppe. Denn das System besitzt eine Einsklasse, bestehend aus den Matrizenringen über \(P\), also den \(P\) ähnlichen Algebren. Es besitzt ferner zu jeder Klasse \((A)\) die inverse \((\overline A)=(A)^{-1}\), bestehend aus den zu \(\overline A\) ähnlichen Algebren, wo \(\overline A\) zu \(A\) reziprok isomorph ist. Das folgt aus dem Vertauschungssatz §5, 3. vermöge der Spezialisierung \(S=A\). Denn \(A\) läßt sich irreduzibel in \(A\) selbst einbetten; es kommt also \(B=P\) und \(A_A=P_f\).\footnote{Die feineren Sätze über die Gruppe der Algebrenklassen benutzen die Theorie der Faktorensysteme; darauf soll an dieser Stelle nicht eingegangen werden. Es sei bemerkt, daß bis zu dieser Stelle keine Überlegungen über kommutative Körper auftreten; z. B. wurde die Tatsache, daß der Rang \((A:P)\) eine Quadratzahl ist, weder benutzt noch bewiesen.}

\paragraph{2. Zerfällungskörper einer Algebrenklasse.}
Die Theorie der Zerfällungskörper beruht wieder vollständig auf dem zu Anfang von §5 ausgesprochenen Prinzip; die kommutativen Erweiterungskörper werden in \(A\) beziehungsweise \(\overline A\) dargestellt. Vorausgeschickt sei der

\paragraph{Hilfssatz.}
Ist \(A\) Körper mit dem Zentrum \(P\), ist \(Z\) endlicher kommutativer Erweiterungskörper von \(P\), so wird \(A_Z\) einfach mit Zentrum \(Z\). Denn das gilt nach §4, 1. für das zu \(A_Z\) reziprok isomorphe System \(Z_{\overline A}\). Für jede Klasse \((A)\) von zu \(A\) ähnlichen Algebren über \(P\) erzeugt also \(Z\) eine Erweiterungsklasse \((A)_Z\) von zu \(A_Z\) ähnlichen einfachen Algebren über \(Z\).

Die zugeordnete Divisionsalgebra \(D\) einer Erweiterungsklasse ergibt sich sofort aus dem Vertauschungssatz aus §5, 3.:

\paragraph{Satz über die zugeordnete Divisionsalgebra einer Erweiterungsklasse.}
Es sei \((A)_Z\) Erweiterungsklasse, \(Z\) bedeute eine irreduzible Einbettung, also Darstellung, von \(Z\) in \(A_f\). Dann wird
\[
        (A)_Z=(D),
\]
wo \(D\) aus der Gesamtheit der mit \(Z\) elementweise vertauschbaren Elemente aus \(A_f\) besteht. Man hat nur in §5, 3. das dortige einfache System \(S\) durch \(Z\) zu ersetzen unter Berücksichtigung, daß \(Z_A\) und \(A_Z\) reziprok isomorph sind. Handelt es sich um reduzible Einbettung von \(Z\), so ergibt die Gesamtheit der mit \(Z\) elementweise vertauschbaren Elemente einen Matrizenring über \(D\).

Ein kommutativer Erweiterungskörper \(Z\) von \(P\) heißt bekanntlich Zerfällungskörper der Klasse \((A)\), wenn die Erweiterungsklasse \((A)_Z\) vollständig zerfällt, d. h. gleich der Einsklasse über \(Z\) wird, also der Klasse aller Matrizenringe über \(Z\).

\paragraph{Satz über die Charakterisierung der Zerfällungskörper.}
Ein endlicher kommutativer Erweiterungskörper \(Z\) von \(P\) ist dann und nur dann Zerfällungskörper der Klasse \((A)\), wenn seine irreduzible Einbettung in \(A_f\) einen maximalen kommutativen Unterkörper von \(A_f\) ergibt.

Denn die Eigenschaft von \(Z\), Zerfällungskörper von \((A)\) zu sein, ist gleichbedeutend mit \((D)=(Z)\). Es muß also nach dem Satz über die zugeordnete Divisionsalgebra die irreduzible Einbettung \(Z\) maximal kommutativer Unterkörper sein. Ist diese Bedingung erfüllt, so wird notwendig \(D=Z\), da jedes \(a\) aus \(D\) durch Adjunktion zu \(Z\) einen kommutativen Erweiterungskörper erzeugt.

\paragraph{Bemerkungen.}
1. Es folgt zugleich, daß ein maximal-kommutativer Unterkörper \(Z\) bei irreduzibler Einbettung einen maximal-kommutativen Unterring erzeugt. Denn die Gesamtheit der mit \(Z\) elementweise vertauschbaren Elemente bildet einen Körper.

2. Der Satz ergibt zugleich die Existenz der Zerfällungskörper; denn die zugeordnete Divisionsalgebra \(A\) besitzt sicher maximal-kommutative Unterkörper.

\paragraph{3. Rangrelationen, Index, unendliche kommutative Erweiterungen.}
Die Rangaussage aus §5, 4. ergibt zunächst die bekannte Tatsache, daß der Rang einer einfachen Algebra nach dem Zentrum eine Quadratzahl ist. Denn ist \(Z\) maximal-kommutativ in \(A\), so kommt, da in \(A\) jede Einbettung irreduzibel ist,
\[
        (Z:P)(Z:P)=(A:P),
\]
somit
\[
        (A:P)=m^2,\qquad (Z:P)=m,
\]
wo \(m\) der Schursche Index ist, der also zugleich den Grad aller in der Divisionsalgebra \(A\) selbst einbettbaren Zerfällungskörper bedeutet. Weiter kommt für den Grad \(n\) eines allgemeinen Zerfällungskörpers
\[
        n=mr,
\]
etwa wegen \((A_r:P)=m^2r^2\) oder auch nach der Rangrelation §4, 3., letzteres, weil der Index \(m\) auch gleich der Anzahl der einfachen Komponenten von \(A_Z\) ist, das Matrizenring über \(Z\) vom Rang \(m^2\) wird. Ist allgemeiner \(A_Z\simeq D_l\), ist \(L\) die irreduzible Einbettung von \(A\) in \(A_r\), und \(d^2\) der Rang \((D:L)\) der Divisionsalgebra \(D\) über ihrem Zentrum \(Z\), so kommt
\[
        l\,d=mr .
\]
Denn nach §5, 4. und dem Satz über die zugeordnete Divisionsalgebra aus 2. kommt
\[
        (A_r:P)=(L:P)(D:P),
\]
also
\[
        m^2r^2=l^2d^2 .
\]

Aus der Existenz der Zerfällungskörper folgen für unendliche kommutative, algebraische oder transzendente, Erweiterungskörper \(\Omega\) von \(P\) die Tatsachen: die Erweiterungsklasse \((A)_\Omega\) ist eine Klasse einfacher Algebren mit dem Zentrum \(\Omega\). Ist insbesondere \(\Omega\) algebraisch abgeschlossen, so wird \(A_\Omega\) Matrizenring vom Grad \(m\). Der Index ist also gleich der „absoluten Komponentenzahl“. Denn ist \(Z\) Zerfällungskörper von \(A\) und \(\Omega'\) das Kompositum von \(\Omega\) und \(Z\), so wird \(A_{\Omega'}\) Matrizenring über \(\Omega'\), also ohne Radikal und mit Zentrum \(\Omega'\). Somit ist auch \(A_\Omega\) ohne Radikal und sein Zentrum \(\Omega\), da ein von \(\Omega\) verschiedenes Element von \(A_\Omega\) auch im Zentrum von \(A_{\Omega'}\) läge, also Körper. \(A_\Omega\) ist also einfache Algebra über \(\Omega\), sicher Matrizenring über \(\Omega\), wenn \(\Omega\) algebraisch abgeschlossen ist.\footnote{Im allgemeinen gibt es auch „minimale“ Zerfällungskörper, d. h. solche, für die kein echter Unterkörper Zerfällungskörper ist, von allen Gradzahlen. Vgl. Brauer--Noether, die in Anm. 2) zitierte Arbeit.}

\paragraph{4. Existenz von separablen Zerfällungskörpern.}
Jede Klasse \((A)\) besitzt separable Zerfällungskörper, sogar solche, die in \(A\) selbst einbettbar sind.\footnote{Vgl. G. Köthe, Über Schiefkörper mit Unterkörpern zweiter Art über dem Zentrum, Journ. f. Math. 166 (1932), S. 182--184. Der vorliegende viel einfachere Beweis geht auf eine Bemerkung von M. Zorn zurück.}

Es sei der Grundkörper \(P\) von der Charakteristik \(p\), der Index
\[
        m=s p^f
\]
mit \(s\) zu \(p\) teilerfremd. Ist dann \(Z\) ein Zerfällungskörper von \(A\) vom Grade \(m\) und \(Z_1\) die darin enthaltene Erweiterung erster Art, also
\[
        (Z:Z_1)=p^f,
\]
so wird der Index von \(A_{Z_1}\sim D\) gleich \(p^f\). Es soll die Existenz eines separablen Erweiterungskörpers \(Z_2\) von \(Z_1\) nachgewiesen werden, derart, daß der Index von \(D_{Z_2}\) ein echter Teiler von \(p^f\) wird. Nun ist \(D_\Omega\) mit algebraisch-abgeschlossenem \(\Omega\) nach 3. ohne Radikal; die reduzierte Diskriminante von \(D\) verschwindet also nicht. Es gibt also mindestens ein Element \(d\) aus \(D\) mit nichtverschwindender reduzierter Spur; dabei kann \(d\) nicht schon im Grundkörper \(Z_1\) liegen, da die Spur von jedem Element \(\alpha\) aus \(Z_1\) gleich \(p^f\alpha\) wird, also verschwindet. \(Z_2=Z_1(d)\) wird also echter, und zwar separabler Erweiterungskörper; der Index von \(D_{Z_2}\) wird echter Teiler von \(p^f\). Die endlich oftmalige Wiederholung führt zu einem separablen Zerfällungskörper von \((A)\), dessen Grad \(m\) mit dem Index von \(A\) übereinstimmt.

\subsection*{§8. Zerfällungskörper, Abspaltungskörper und galoissche Theorie im Kommutativen}

Um von den Zerfällungskörpern der über ihrem Zentrum einfachen Systeme zu denjenigen beliebiger einfacher Systeme überzugehen, ist noch die Zerfällungstheorie des Zentrums zu entwickeln, der noch eine zu §6, 3. parallel laufende galoissche Theorie hinzugefügt werden soll. Alle in diesem Paragraphen auftretenden Körper und Systeme sind kommutativ.

\paragraph{1. Zerfällungs- und Abspaltungskörper kommutativer einfacher Systeme.}
Das kommutative System \(Z\) sei einfach über \(P\), also Körper, \(\Omega\) ein algebraisch abgeschlossener Erweiterungskörper von \(P\). Dann gilt bekanntlich: dann und nur dann bleibt \(Z_\Omega\) System ohne Radikal, wenn \(Z\) separabel, also Erweiterung erster Art, über \(P\) ist. Denn dann und nur dann besitzt es so viele isomorphe Abbildungen ersten Grades in \(\Omega\), wie sein Rang nach \(P\) beträgt, also auch ebenso viele verschiedene Darstellungen, die hier mit Darstellungsklassen übereinstimmen.

Die Tatsache, daß in \(Z_\Omega\) ein Radikal auftreten kann, also nicht notwendig direkte Summenzerlegung in absolut einfache Komponenten statthat, führt zu den folgenden Definitionen: ein Erweiterungskörper \(\Lambda\) von \(P\) heißt Zerfällungskörper, wenn \(Z_\Lambda\) in Kompositionsfaktoren ersten Grades zerfällt; mit anderen Worten, wenn alle absolut irreduziblen Darstellungen schon in \(\Lambda\) liegen. Ein Erweiterungskörper \(T\) von \(P\) heißt Abspaltungskörper, wenn \(Z_T\) mindestens einen Kompositionsfaktor ersten Grades abspaltet; mit anderen Worten, wenn mindestens eine absolut irreduzible Darstellung von \(Z\) in \(T\) existiert.

Zerfällungs- und Abspaltungskörper lassen sich charakterisieren nach dem

\paragraph{Satz.}
Ein Körper \(T\) ist dann und nur dann Abspaltungskörper, wenn er einen zu \(Z\) isomorphen Unterkörper \(Z'\) enthält; ein Körper \(\Lambda\) ist dann und nur dann Zerfällungskörper, wenn er einen Unterkörper \(\Gamma\) enthält, der dem zu \(Z\) gehörigen galoisschen Körper isomorph wird.

Das folgt direkt aus der Definition von Zerfällungs- und Abspaltungskörper vermöge absolut irreduzibler Darstellungen. Insbesondere ergibt sich in Analogie zum Nichtkommutativen die Folgerung: ein Körper \(Z'\) ist dann und nur dann minimaler Abspaltungskörper, d. h. kein echter Unterkörper ist Abspaltungskörper, wenn er zu \(Z\) isomorph ist. Ein minimaler Abspaltungskörper ist dann und nur dann zugleich Zerfällungskörper, wenn \(Z\) normal, also galoisscher Körper über \(P\) ist.

Hierin liegt der Grund dafür, eine einfache Algebra normal über ihrem Zentrum zu nennen. Die Tatsache, daß es innerhalb eines festen algebraisch abgeschlossenen \(\Omega\) über \(P\) nur einen einzigen minimalen Zerfällungskörper gibt, den zu \(Z\) gehörigen galoisschen, im Gegensatz zu den im allgemeinen unendlich vielen nichtisomorphen im Nichtkommutativen, hat ihren Grund in der eindeutigen direkten Summenzerlegung im Kommutativen im Gegensatz zu der nur bis Operatorisomorphie eindeutigen im Nichtkommutativen; oder, was auf dasselbe hinauskommt, in den nur endlich vielen verschiedenen absolut irreduziblen Darstellungen an Stelle der unendlich vielen verschiedenen im Nichtkommutativen, die allerdings derselben Klasse angehören.

\paragraph{2. Hyperkomplexe Begründung der galoisschen Theorie im Fall separabler Körper \(Z/P\).}
In genauer Analogie zu §6, 3. gilt hier die Theorie der Isomorphismen für beliebige, über \(P\) separable \(Z\); im Fall eines galoisschen \(Z\) läßt sich dann der Übergang zur üblichen Automorphismengruppe machen.

\paragraph{Voraussetzungen.}
Das einfache System \(Z\) über \(P\) sei endlicher separabler Erweiterungskörper, \(\Omega\) bedeute einen über \(P\) endlichen oder unendlichen Zerfällungskörper, \(Z'=Z^{(1)}\) einen zu \(Z\) isomorphen Unterkörper, \(\Gamma\) den zugehörigen galoisschen. Nach 1. gilt also die direkte Summenzerlegung
\[
        Z_\Omega=r^{(1)}+\cdots+r^{(n)}
        =e^{(1)}Z_\Omega+\cdots+e^{(n)}Z_\Omega
        =e^{(1)}\Omega+\cdots+e^{(n)}\Omega .
\]
Die durch \(e^{(i)}\) erzeugte Darstellung \(Z\to Z^{(i)}\) mit \(Z^{(i)}\subseteq\Gamma\) ist also definiert durch
\[
        e^{(i)}z=e^{(i)}z^{(i)}.
\]

\paragraph{Hilfssatz 1.}
Die \(n\) durch \(e^{(1)},\ldots,e^{(n)}\) erzeugten Darstellungen sind alle verschieden. Der Beweis ist wie in §6, 3., oder auch nach 1.; denn sie gehören verschiedenen Darstellungsklassen an.

\paragraph{Hilfssatz 2.}
Ist \(T\) Zwischenkörper von \(P\) und \(Z\), ist \(s\) der Rang von \(Z\) nach \(T\), so läßt jeder Isomorphismus von \(T\) in \(\Omega\) mindestens, und daher genau, \(s\) Fortsetzungen zu.

Der Beweis ist wie in §6, 3. Daß hier „mindestens“ zu „genau“ verschärft wird, folgt aus der eindeutigen Zerlegung von \(Z_\Omega\), nach der \(Z\) nicht mindestens, sondern genau \(n\) Isomorphismen in \(\Omega\) besitzt.

Wie in §6, 3. definiert man jetzt die durch \(T\) induzierte Klasseneinteilung \(\mathfrak I_T\) der endlichen Menge \(\mathfrak I\) der Isomorphismen von \(Z\): alle und nur die Isomorphismen aus \(\mathfrak I\) werden in \(\mathfrak I_T\) als äquivalent betrachtet, die auf \(T\) denselben Isomorphismus induzieren. Man beweist den

\paragraph{Hauptsatz.}
Ist \(\mathfrak I_T\) die durch \(T\) induzierte Klasseneinteilung, so ist \(T\) maximaler Unterkörper gegenüber dieser Klasseneinteilung.

Von hier aus kann man für den Fall eines galoisschen \(Z\) genau wie in §6, 3. durch Zusammensetzung mit dem Reziproken eines Isomorphismus zur Automorphismengruppe und zur üblichen Fassung übergehen. Die entsprechende Aussage für die Untergruppen ergibt sich aus der Betrachtung der Idempotente, die an sich von Interesse ist.

\paragraph{3. Die Idempotente von \(Z_\Omega\).}
Es durchlaufe \(S\) die \(n\) Substitutionen, die \(Z\) in die einzelnen \(Z^{(i)}\) überführen, also die Zusammensetzung der Isomorphismen \(Z\to Z_1\) und \(Z\to Z^{(i)}\), wobei der erstere den zu \(Z\to Z_1\) reziproken bedeutet. \(Z\) braucht nicht galoissch zu sein. Für die Elemente \(a\) aus \(Z_\Omega\) seien die Substitutionen \(S\) als Operatoren definiert durch die Festsetzung, daß sie auf \(Z\) die Identität induzieren; zu jedem \(a\) ist somit das in \(Z_{(i)}\) liegende konjugierte \(a^S\) definiert. Bei diesen Festsetzungen gilt der

\paragraph{Satz.}
Die \(n\) Idempotente \(e^{(i)}\) von \(Z_\Omega\) sind konjugiert:
\[
        e^{(i)}=e^S,\qquad e=e^{(1)};
\]
sie liegen in den \(n\) konjugierten Erweiterungssystemen \(Z\Omega\).

Denn die Anwendung von \(S\) führt die Definitionsrelationen
\[
        ez=ez^{(1)},\qquad e^2=e
\]
über in
\[
        e^S z=e^S z^{(i)},\qquad (e^S)^2=e^S .
\]
Wegen der Eindeutigkeit der Zerlegung sind auch die Idempotente eindeutig bestimmt; also \(e^S=e^{(i)}\). Da ferner \(Z\) Abspaltungskörper ist, also die Darstellung \(ez=ez^{(1)}\) schon vermittelt, liegt \(e\) schon in \(Z\Omega\), und damit \(e^S\) in \(Z^S\Omega\).

Ist speziell \(Z\) galoissch, so folgt unmittelbar der auf die Untergruppen bezügliche Teil des Hauptsatzes der galoisschen Theorie; denn ist \(\mathfrak H\) Untergruppe der galoisschen Gruppe \(\Gcal\), so gestattet, wegen der eindeutigen Bestimmtheit der Komponenten einer direkten Summe, das Element
\[
        \sum_{H\in\mathfrak H} e^H
\]
alle Substitutionen aus \(\mathfrak H\), aber keine davon verschiedenen. Da die Gruppe \(\Gcal\) für \(Z\) definiert war, handelt es sich um die galoissche Theorie von \(Z\); wegen der Isomorphie ist damit die für \(Z'\) mit erledigt.

\paragraph{4. Zusammenhang der Idempotente mit den komplementären Basen.}
\(Z\) wird wieder nicht notwendig als galoissch vorausgesetzt.

\paragraph{Satz.}
Ist \(a_1,\ldots,a_n\) eine Basis von \(Z\) über \(P\), wird das Idempotent
\[
        e=a_1\beta_1+\cdots+a_n\beta_n
\]
gesetzt, also
\[
        e^S=a_1\beta_1^S+\cdots+a_n\beta_n^S,
\]
so bedeutet
\[
        \beta_1^S,\ldots,\beta_n^S
\]
das Bild der Komplementärbasis \(b_1,\ldots,b_n\) zu \(a_1,\ldots,a_n\) in der durch \(e^S\) vermittelten Abbildung.

Bei der durch \(e^S\) vermittelten Abbildung \(e^Sx=e^Sx^S\), insbesondere \(e^Sa_i=e^Sa_i^S\), gelten wegen
\[
        e^Se^S=e^S,\qquad e^Se^R=0\quad(S\ne R)
\]
die Zuordnungen \(e^S\mapsto1\), \(e^R\mapsto0\). Also kommt
\[
        1=a_1^S\beta_1^S+\cdots+a_n^S\beta_n^S,\qquad
        0=a_1^R\beta_1^S+\cdots+a_n^R\beta_n^S\quad(S\ne R).
\]
In Matrizen geschrieben ist dies die Definitionsgleichung der komplementären Basen:
\[
\begin{pmatrix}
a_1^1 & \cdots & a_n^1\\
\vdots & \ddots & \vdots\\
a_1^n & \cdots & a_n^n
\end{pmatrix}
\begin{pmatrix}
\beta_1^S\\ \vdots\\ \beta_n^S
\end{pmatrix}
=
\begin{pmatrix}
0\\ \vdots\\ 1\\ \vdots\\ 0
\end{pmatrix}.
\]
Der Übergang zur Spurendefinition folgt daraus, daß für
\[
        a_i=\sum_S a_i^S e^S,\qquad b_i=\sum_S b_i^S e^S
\]
die Matrixrelation nach Vertauschung der Zeilen übergeht in
\[
        \delta_{ij}=\Sp(a_i b_j),
\]
wobei die \(a_i\) und \(b_i\) in \(Z\) liegen, da sie alle Substitutionen \(S\) gestatten.\footnote{Vgl. Darstellungstheorie §25.}\footnote{Beziehungen zwischen den komplementären Basen und den Idempotenten finden sich erstmalig bei Dedekind, Zur Theorie der aus \(n\) Haupteinheiten gebildeten komplexen Größen; vgl. Bd. II der Gesammelten Werke, S. 4--5.}

Die Kontragredienz der komplementären Basen folgt hier aus der Tatsache, daß die Idempotente \(e^S\) Invarianten von \(Z\) sind, unabhängig von der zugrunde gelegten Basis, daß also alle
\[
        a_1\beta_1^S+\cdots+a_n\beta_n^S
\]
in sich übergehen beim Übergang zu einer anderen Basis
\[
        a_1',\ldots,a_n'
\]
von \(Z/P\).

\subsection*{§9. Zerfällungskörper und Abspaltungskörper beliebiger Systeme}

\paragraph{1. Definitionen und Reduktion auf den einfachen Fall.}
Den in §8 gegebenen Definitionen entsprechen im allgemeinen Fall die Definitionen: sei \(S\) hyperkomplex über \(P\). Ein kommutativer Erweiterungskörper \(\Lambda\) von \(P\) heißt Zerfällungskörper von \(S\), wenn \(S_\Lambda\) bei einer Kompositionsreihe aus einseitigen Idealen in absolut einfache Kompositionsfaktoren zerfällt; mit anderen Worten, wenn alle irreduziblen Darstellungen von \(S\) in \(\Lambda\) schon absolut irreduzibel sind. Ein Erweiterungskörper \(T\) von \(P\) heißt Abspaltungskörper, wenn \(S_T\) mindestens einen absolut einfachen Kompositionsfaktor abspaltet, wenn also mindestens eine in \(T\) irreduzible Darstellung schon absolut irreduzibel ist.

Diese Definitionen enthalten offenbar die in §8 für das Kommutative, in §7 für einfache Algebren gegebenen in sich; das letztere, weil hier \(A_\Lambda\) bei jeder kommutativen Erweiterung des Zentrums vollständig reduzibel ist und somit Kompositionsfaktoren und Komponenten der direkten Summenzerlegung übereinstimmen. Volle Matrizenringe über dem Zentrum und nur diese ergeben aber hier absolut einfache Komponenten, also auch absolut irreduzible Darstellungen. Da ferner \(A_\Lambda\) zweiseitig einfach ist, also in operatorisomorphe Komponenten zerfällt, so ergibt das Abspalten einer absolut einfachen Komponente schon das vollständige Zerfallen: Abspaltungs- und Zerfällungskörper fallen zusammen.

Aus den Definitionen ergeben sich unmittelbar die Tatsachen: die Zerfällungskörper von \(S\) sind gegeben als Vereinigungskörper je eines Zerfällungskörpers der in \(P\) irreduziblen Darstellungen; Abspaltungskörper von \(S\) ist jeder Abspaltungskörper einer in \(P\) irreduziblen Darstellung. Wegen des eineindeutigen Entsprechens der einfachen Systeme, Komponenten des Restklassenrings nach dem Radikal, und der in \(P\) irreduziblen Darstellungen genügt also die Betrachtung einfacher Systeme. Hier werden die Resultate sich durch Kombination von §7 und §8 ergeben.

\paragraph{2. Zerfällungskörper und Abspaltungskörper einfacher Systeme.}
Wir betrachten vorerst den Fall eines einfachen Systems \(A\) mit separablem Zentrum \(Z\) über \(P\). Aus §8, 1. und §8, 3. kommt: die der direkten Summenzerlegung von \(Z_\Omega\) entsprechende zweiseitige Zerlegung von \(A_\Omega\), wo \(\Omega\) einen Zerfällungskörper von \(Z\) bedeutet, ergibt \(n\) konjugierte, zu \(A\) isomorphe Komponenten
\[
        e^S A_\Omega
\]
von \(A_\Omega\), die einfache Algebren über ihrem Zentrum
\[
        e^S Z_\Omega=e^S Z
\]
werden. Dabei sind die Substitutionen \(S\) als Operatoren für \(A_\Omega\) definiert durch die Festsetzung, daß sie auf \(A\) die Identität induzieren.

Denn nach §8, 3. sind die Idempotente konjugiert; also gilt das gleiche von den Komponenten \(e^S A_\Omega\) nach Definition der Substitutionen für \(A\). Da \(A\) zweiseitig einfach ist, wird \(e^S A_\Omega\) ringisomorph zu \(A\). Dabei wird \(e^SA_\Omega=e^S A_{Z^{S}}\), wegen \(e^S Z_\Omega=e^S Z\); das Zentrum fällt also mit dem Koeffizientenbereich zusammen, und es handelt sich bei den Komponenten um normale Algebren.

Die Resultate von §7, 2. ergeben jetzt weiter: ist \(A\) einfaches System über \(P\) mit separablem Zentrum \(Z\) über \(P\), so wird auch \(A_\Omega\) ohne Radikal, also vollständig reduzibel; dabei kann \(\Omega\) irgendeinen endlichen oder unendlichen Erweiterungskörper bedeuten. Denn nach §7, 2. bleibt \(e^S A_\Omega\) einfach bei jeder Koeffizientenerweiterung, also ohne Radikal. Somit gilt das gleiche für die direkte Summe; diese ist aber \(A_\Omega\), oder entsteht, wenn \(\Omega\) kein Zerfällungskörper von \(Z\) ist, durch Koeffizientenerweiterung aus \(A_\Omega\).

Weiter folgt aus §7, 2. die

\paragraph{Charakterisierung der Abspaltungskörper und Zerfällungskörper.}
Ist \(A\) Divisionsalgebra mit separablem Zentrum \(Z\), so sind die irreduziblen Einbettungen der Abspaltungskörper gegeben durch alle und nur die \(Z\) umfassenden maximalen kommutativen Unterkörper der \(A_f\); Zerfällungskörper sind alle und nur die Vereinigungskörper eines Abspaltungskörpers mit einem Zerfällungskörper des Zentrums, insbesondere mit dem zum Zentrum gehörigen galoisschen Körper. Ein minimaler Abspaltungskörper kann Zerfällungskörper sein, auch ohne daß das Zentrum galoissch ist.

Denn ein Abspaltungskörper muß einen zu \(Z\) isomorphen Körper umfassen; die Einbettungsaussage über die Abspaltungskörper folgt jetzt aus den vorangehenden Tatsachen und aus §7, 2. Ein Zerfällungskörper muß ferner einen Zerfällungskörper \(\Omega\) zu \(Z\) umfassen. Jeder Abspaltungskörper über \(\Omega\) ist aber zugleich Zerfällungskörper, da die Komponenten \(e^S A_\Omega\) einfach sind und ein zu \(\Omega\) isomorphes Zentrum besitzen. Ist \(Z\) galoissch, so fallen also minimale Abspaltungs- und Zerfällungskörper zusammen. Aber auch im nichtgaloisschen Fall kann es, anders als im Kommutativen, minimale Abspaltungskörper geben, die zugleich Zerfällungskörper sind, nämlich sobald ein solcher minimaler Abspaltungskörper den zu \(Z\) gehörigen galoisschen Körper umfaßt. Beispiel: Quaternionenkörper mit zu \(P(\sqrt2)\) isomorphem Zentrum, \(P\) Körper der rationalen Zahlen; der zu \(P(\sqrt2)\) gehörige galoissche Körper ist minimaler Abspaltungs- und Zerfällungskörper.

Handelt es sich schließlich um inseparables Zentrum, so wird \(Z_\Omega\) und damit \(A_\Omega\) System mit Radikal. Sei etwa \(\mathfrak C\) Radikal von \(Z_\Omega\), und \(\mathfrak C A_\Omega\) das Erweiterungsideal in \(A_\Omega\). Dann gilt für \(Z_\Omega/\mathfrak C\) eine direkte Summenzerlegung in konjugierte Komponenten, in genauer Analogie zu §8, 2.; das ergibt genau wie oben die direkte Summenzerlegung von \(A_\Omega/\mathfrak C A_\Omega\), derart, daß die Komponenten \(e^SA_\Omega\) zu \(A\) isomorph sind, mit Zentrum \(e^SZ^\Omega\). Also bleiben auch die Sätze über Abspaltungs- und Zerfällungskörper erhalten. Weiter zeigt die Rangabzählung, daß die einer Kompositionsreihe von \(\mathfrak C A_\Omega\) entsprechenden Kompositionsfaktoren zu \(A_\Omega/\mathfrak C A_\Omega\) operatorisomorph, nicht nur homomorph werden.

Für irreduzible Darstellungen ausgesprochen, ergibt sich also die Zusammenfassung: liegt eine in \(P\) irreduzible Darstellung eines hyperkomplexen Systems vor, so wird die Darstellung dann und nur dann absolut vollständig reduzibel, wenn das zugehörige Zentrum über \(P\) separabel ist. In jedem Fall lassen sich Zerfällungs- und Abspaltungskörper der Darstellung durch Einbettung in das zugehörige einfache System wie oben charakterisieren. Insbesondere wird der niedrigste Grad eines Abspaltungskörpers gleich dem Produkt aus dem Grad des Zentrums mit dem Index der zugehörigen Algebrenklasse über dem Zentrum; der Grad von jedem Abspaltungskörper ist ein Vielfaches davon. Die Darstellung zerfällt, im Sinn der Kompositionsreihen, in so viele verschiedene, aber konjugierte absolut irreduzible Darstellungsklassen, wie der Grad des größten im Zentrum enthaltenen separablen Körpers beträgt. Jede Klasse tritt so oft auf, wie das Produkt aus Index und Grad des Zentrums nach dieser separablen Erweiterung ausmacht.

Für vollkommenen Grundkörper, wo es nur separable Erweiterungskörper gibt, kommen wir so auf bekannte Resultate von I. Schur zurück, unter Zufügung der Einbettungscharakterisierung, die sich schon bei R. Brauer findet.

\begin{center}
(Eingegangen am 8. Juni 1932.)
\end{center}

\end{document}
